\begin{abstract}
    Chemical plants generate high-dimensional sensor streams in which missing values in critical tags can degrade monitoring, control, and operational decision-making. In practice, missing targets are often handled by simple substitution, nearest-neighbor/regression imputers, or generic reconstruction models, frequently under evaluation protocols that leak preprocessing information or fail to reflect deployment conditions. We address this gap with a protocol-first, leakage-safe study for reconstructing a critical signal in an industrial sulfate-process unit under extreme target sparsity, using a \emph{target-conditional denoising autoencoder} (CDAE) aligned with the single-sensor outage setting (target unavailable at inference time while auxiliary variables remain observed). Phase~1 applies split-first, train-only preprocessing (feature selection, imputation, and scaling), persists all artifacts for byte-level reproducibility, and achieves strong held-out performance ($R^2\!\approx\!\phaseOneRtwo$). Phase~2 performs mechanism-aware controlled evaluation by masking only the \emph{test} target and scoring only masked indices under missing completely at random (MCAR), missing at random (MAR), and missing not at random (MNAR), using the same blocked temporal split. With 80\% target masking, the CDAE attains $R^2=\mcarRtwo$ (MCAR), $R^2=\marRtwo$ (MAR), and $R^2=\mnarRtwo$ (MNAR), with $n=\mcarN$ masked samples per mechanism and MAE between 2.3 and 2.5~EU on a 0--86~EU scale. The degradation under MNAR is consistent with value-dependent missingness under temporal distribution shift. Residual diagnostics indicate non-Gaussian, heteroscedastic errors concentrated in the low-flow, low-density regime. The main contribution is a deployment-aligned, leakage-safe evaluation framework for industrial virtual sensing under extreme sparsity, complemented by evidence that target conditioning drives global temporal accuracy while mask weighting ($\lambda{=}1$) reduces low-flow bias at a trade-off in global MAE.
    \end{abstract}